\chapter{Transformer 推理显微镜：逐层张量变化与 RoPE 实测}

\section{Lab 04：大模型推理的微观全貌（零起点极细拆解）}

如果你对「大模型推理到底在算什么」、「张量每一步怎么变」感到模糊，本模块就是为你量身打造的\textbf{显微镜}。

无论模型是 0.5B、7B 还是 70B，无论叫 Llama-3、Qwen2.5 还是 DeepSeek，其自回归推理的\textbf{微观数学本质完全一致}。

\vspace{0.8em}\hrule\vspace{0.8em}

\section{[指南] 1. 宏观全貌：从一句话到下一个字的 8 个步骤}

假设你对模型说：\texttt{"人工智能"}，模型要推理出下一个字 \texttt{"的"}。整个过程分为严格的 8 步：

\begin{figure}[htbp]
\centering
\begin{tikzpicture}[
  stage/.style={draw=primary!80, fill=primary!5, rounded corners=1.5mm, align=center, font=\small\sffamily, inner sep=4.5pt, line width=0.8pt, minimum width=11cm},
  arr/.style={-{Stealth[scale=0.9]}, draw=primary, line width=1pt}
]
\node[stage] (s0) {原始输入文本：\texttt{"人工智能"}};
\node[stage, below=0.4cm of s0] (s1) {1. Tokenizer 分词 $\rightarrow$ Token IDs: \texttt{[102, 305]} \quad (Shape: $[2]$)};
\node[stage, below=0.4cm of s1] (s2) {2. Embedding 查表 $\rightarrow$ 稠密向量矩阵 \quad (Shape: $[2, d_{\mathrm{model}}]$)};
\node[stage, below=0.4cm of s2] (s3) {3. RMSNorm 归一化 $\rightarrow$ 消除激活值方差偏置 \quad (Shape: $[2, d_{\mathrm{model}}]$)};
\node[stage, below=0.4cm of s3] (s4) {4. 多头自注意力 (Self-Attention + RoPE 旋转位置编码) \quad (Shape: $[2, d_{\mathrm{model}}]$)};
\node[stage, below=0.4cm of s4] (s5) {5. 残差连接 + SwiGLU / MLP 非线性知识映射 \quad (Shape: $[2, d_{\mathrm{model}}]$)};
\node[stage, below=0.4cm of s5, fill=secondary!10, draw=secondary] (s6) {6. 循环堆叠 $L$ 个 Transformer 块 (重复步骤 3 $\sim$ 5)};
\node[stage, below=0.4cm of s6] (s7) {7. Final RMSNorm + LM Head 词表投影 $\rightarrow$ Logits \quad (Shape: $[2, V]$)};
\node[stage, below=0.4cm of s7, fill=accent!10, draw=accent] (s8) {8. 采样算法 (Top-P / Temperature) $\rightarrow$ 生成下一个 Token ID: 520 (\texttt{"的"})};

\draw[arr] (s0) -- (s1);
\draw[arr] (s1) -- (s2);
\draw[arr] (s2) -- (s3);
\draw[arr] (s3) -- (s4);
\draw[arr] (s4) -- (s5);
\draw[arr] (s5) -- (s6);
\draw[arr] (s6) -- (s7);
\draw[arr] (s7) -- (s8);
\end{tikzpicture}
\caption{Transformer 单步推理张量生命周期微观全貌}
\end{figure}

\vspace{0.8em}\hrule\vspace{0.8em}

\section{[观察] 2. 张量形状流（Tensor Shape Flow Tracker）}

理解大模型推理的关键，是死死盯住\textbf{每一步张量的形状（Shape）}。
下表用一个\textbf{假设的}微型配置来说明形状如何流动（注意：脚本 \texttt{01\_tensor\_trace\_step\_by\_step.py} 为了能把每个数字打印出来，用的是更小的配置：$d_{\text{model}} = 16$、2 个头、词表 11 个字、只有 1 层；下面的流程图描述的是真实模型的 $L$ 层结构）。

假设微型模型配置：
\begin{itemize}
  \item 序列长度 $T = 3$（例如输入 3 个词）
  \item 隐藏层维度 $d_{\text{model}} = 64$
  \item 注意力头数 $n_{\text{heads}} = 4$，每个头的维度 $d_{\text{head}} = 64 / 4 = 16$
  \item 词表大小 $\text{vocab\_size} = 1000$
\end{itemize}

\begin{center}
\small
\begin{tabularx}{\textwidth}{l X X X X }
\toprule
\textbf{处理阶段} & \textbf{算子名称} & \textbf{输入形状 (Input Shape)} & \textbf{输出形状 (Output Shape)} & \textbf{物理含义解释} \\
\midrule
\textbf{分词} & Tokenizer & 字符串文本 & \texttt{[3]} & 文本转整数 ID \\
\textbf{词嵌入} & Embedding Lookup & \texttt{[3]} & \texttt{[3, 64]} & 每个词映射成 64 维空间的一个点 \\
\textbf{归一化} & RMSNorm & \texttt{[3, 64]} & \texttt{[3, 64]} & 均方根缩放，保持方差稳定 \\
\textbf{QKV投影} & Linear ($W_q, W_k, W_v$) & \texttt{[3, 64]} & \texttt{[3, 64]} (各 1 个) & 生成查询、键、值向量 \\
\textbf{多头拆分} & Reshape \& Transpose & \texttt{[3, 64]} & \texttt{[4, 3, 16]} & 4 个注意力头并行关注不同维度的特征 \\
\textbf{位置编码} & RoPE (旋转编码) & \texttt{[4, 3, 16]} & \texttt{[4, 3, 16]} & 对 Q 和 K 按位置角频率旋转，注入顺序感 \\
\textbf{注意力打分} & $Q \cdot K^T / \sqrt{d_k}$ & Q:\texttt{[4, 3, 16]}, K:\texttt{[4, 3, 16]} & \texttt{[4, 3, 3]} & 第 $i$ 个词与第 $j$ 个词的关联度矩阵 \\
\textbf{因果掩码} & Causal Mask & \texttt{[4, 3, 3]} & \texttt{[4, 3, 3]} & 将右上角置为 $-\infty$（禁止看到未来的词） \\
\textbf{注意力归一} & Softmax & \texttt{[4, 3, 3]} & \texttt{[4, 3, 3]} & 转换成概率百分比分布（每一行和为 1） \\
\textbf{值加权} & Score $\cdot V$ & Scores:\texttt{[4, 3, 3]}, V:\texttt{[4, 3, 16]} & \texttt{[4, 3, 16]} & 提取汇聚上下文信息 \\
\textbf{多头合并} & Transpose \& Reshape & \texttt{[4, 3, 16]} & \texttt{[3, 64]} & 拼接 4 个头的特征 \\
\textbf{输出投影} & Linear ($W_o$) & \texttt{[3, 64]} & \texttt{[3, 64]} & 特征融合 \\
\textbf{残差连接} & $X + \text{Attn}(X)$ & \texttt{[3, 64]} + \texttt{[3, 64]} & \texttt{[3, 64]} & 保留原始输入，防止梯度/特征退化 \\
\textbf{前馈网络} & SwiGLU / MLP & \texttt{[3, 64]} & \texttt{[3, 64]} & 升维到例如 172 维激活后再降维回 64 维 \\
\textbf{残差连接} & $X + \text{FFN}(X)$ & \texttt{[3, 64]} + \texttt{[3, 64]} & \texttt{[3, 64]} & 第二层残差 \\
\textbf{词表投影} & LM Head ($W_{\text{head}}$) & \texttt{[3, 64]} & \texttt{[3, 1000]} & 将每个位置映射为词表中每个词的得分 \\
\textbf{预测下个词} & Slice Last Position & \texttt{[3, 1000]} & \texttt{[1000]} & \textbf{我们只取最后一个位置（第 3 个词）的输出} \\
\textbf{采样策略} & Greedy / Top-P & \texttt{[1000]} & 整数（标量） & 采样出第 4 个词的 Token ID！ \\
\bottomrule
\end{tabularx}
\end{center}

\vspace{0.8em}\hrule\vspace{0.8em}

\section{[实验] 3. 逐个核心算子的底层细节与数学公式}

\subsection{(1) 为什么注意力机制必须有「因果掩码（Causal Mask）」？}

在文本生成中，我们不能让前面的词“偷看”后面的词。
对于序列长度为 3 的输入：

\[
\text{Mask} = \begin{bmatrix}
0 & -\infty & -\infty \\
0 & 0 & -\infty \\
0 & 0 & 0
\end{bmatrix}
\]

当注意力分数矩阵与该 Mask 相加后，在执行 $\text{Softmax}$ 时：

\[
e^{-\infty} = 0
\]

被掩码的位置权重精确变成 \textbf{0}！这样第 1 个词只能看第 1 个词，第 2 个词只能看第 1 和第 2 个词，保证严格的单向自回归因果性。

\subsection{(2) RMSNorm（均方根归一化）到底算什么？}

现代大模型（LLaMA/Qwen/Mistral）舍弃了传统的 LayerNorm，改用更高效的 RMSNorm：
\begin{itemize}
  \item 不需要计算均值 $\mu$，不需要减均值；
  \item 仅计算向量元素的均方根（Root Mean Square）：
\end{itemize}

\[
\text{RMS}(x) = \sqrt{\frac{1}{d} \sum_{i=1}^{d} x_i^2 + \epsilon}
\]

\begin{itemize}
  \item 归一化公式（$\gamma$ 是可学习的缩放向量）：
\end{itemize}

\[
\bar{x}_i = \frac{x_i}{\text{RMS}(x)} \odot \gamma_i
\]

计算量少了一半，硬件执行速度更快，数值更平稳。

\subsection{(3) SwiGLU（门控前馈网络）算什么？}

现代模型基本不再使用经典的 \texttt{ReLU(W1 x) W2}，而是采用 \textbf{SwiGLU} 结构：
它包含三个权重矩阵：$W_{\text{gate}}, W_{\text{up}}, W_{\text{down}}$：

\[
\text{SwiGLU}(x) = \left( \text{SiLU}(x W_{\text{gate}}) \otimes (x W_{\text{up}}) \right) W_{\text{down}}
\]

\begin{itemize}
  \item $x W_{\text{gate}}$ 通过 $\text{SiLU}$ 激活函数充当“开关门（Gate）”，控制哪些特征信息可以通过。
  \item 与 $x W_{\text{up}}$ 做逐元素乘法（Hadamard product $\otimes$）。
  \item 最终经过 $W_{\text{down}}$ 降维回隐藏层维度。
\end{itemize}

\subsection{(4) 采样：从 Logits 到文字}

最后一步"从 Logits 选出下一个 token"单独放在 \textbf{Lab 07a}，配合第 7 篇学习。

\vspace{0.8em}\hrule\vspace{0.8em}

\section{[执行] 运行显微镜级追踪脚本}

本目录下提供了可以直接运行的 Python 代码：

\begin{enumerate}
  \item \textbf{\textbf{\texttt{01\_tensor\_trace\_step\_by\_step.py}}}：
\end{enumerate}

   \textbf{单步前向传播显微镜}。完整实现了一个微型 Transformer，单步逐行打印所有输入/输出张量 Shape、注意力矩阵数值、Logits 排行榜以及生成循环。
\begin{enumerate}
  \item \textbf{\textbf{\texttt{02\_rope\_and\_attention\_microscope.py}}}：
\end{enumerate}

   \textbf{位置编码与因果掩码显微镜}。无位置编码时的置换等变性、RoPE 的相对位置性质、长程衰减的正确表述（只在 q、k 相关时的期望上成立）、因果掩码。

对应理论：\textbf{第 4 篇}。依赖：纯 Python 标准库。

\vspace{0.8em}\hrule\vspace{0.8em}

\section{[链接] 说明与下一步}

\begin{itemize}
  \item 脚本 01 的 "Decode" 步骤为了看清张量，故意把整段序列重算一遍（没有 KV Cache）；权重是随机的，所以输出没有语义。
  \item 脚本 02 已修正：旧版把"点积交换律"当成了"没有语序感"的证明，并声称单一频率的 RoPE 有长程衰减（实际上单一频率下点积是 $\cos(\Delta\theta)$，是周期函数）。新版演示了真正的\textbf{置换等变性}；长程衰减实验对比了 q、k 完全相关 / 部分相关 / 独立随机三种情况——只有相关时点积的期望才随距离衰减，独立随机时完全不衰减。
  \item 想看一个\textbf{真正训练过}的模型，请继续做 \textbf{Lab 05} 与 \textbf{Lab 07b}。完整顺序见 \textbf{LEARNING\_PATH.md}。
\end{itemize}



\section{实验核心源码精选与解析}

本实验配套有完整可运行的工业级验证代码，重点实现细节如下：


\subsection{源码实现：\texttt{01\_tensor\_trace\_step\_by\_step.py}}

\lstinputlisting[language=Python, caption={\texttt{01\_tensor\_trace\_step\_by\_step.py}}]{code/lab04_transformer_microscope_01_tensor_trace_step_by_step.py}


\subsection{源码实现：\texttt{02\_rope\_and\_attention\_microscope.py}}

\lstinputlisting[language=Python, caption={\texttt{02\_rope\_and\_attention\_microscope.py}}]{code/lab04_transformer_microscope_02_rope_and_attention_microscope.py}

